\newpage
\chapter{KESIMPULAN DAN SARAN} \label{Bab V}

\section{Kesimpulan} \label{V.Kesimpulan}
Berdasarkan hasil penelitian yang telah dilakukan, dapat disimpulkan bahwa penelitian ini berhasil menjawab seluruh rumusan masalah melalui perancangan, pengembangan, serta evaluasi model klasifikasi berbasis multimodal. Proses penelitian dimulai dari pengumpulan dan anotasi data, pembangunan model unimodal dan multimodal, hingga analisis performa menggunakan berbagai skenario eksperimen. Hasil yang diperoleh menunjukkan bahwa pendekatan multimodal mampu memberikan peningkatan performa yang signifikan dibandingkan pendekatan unimodal. Selain itu, penelitian ini juga memberikan \textit{insight} terkait efektivitas metode \textit{fusion} serta kompleksitas arsitektur dalam pemodelan hubungan antar modalitas. Adapun kesimpulan yang diperoleh adalah sebagai berikut:

\begin{enumerate}
    \item Penelitian ini berhasil membangun model klasifikasi biner berbasis multimodal dengan mengintegrasikan CLIP sebagai \textit{encoder} visual dan ELECTRA sebagai \textit{encoder} tekstual melalui pendekatan \textit{intermediate fusion}. Model yang dikembangkan mampu memproses dan menggabungkan informasi visual dan tekstual terintegrasi, sehingga dapat memahami konteks meme secara lebih menyeluruh. Dengan memanfaatkan keunggulan masing-masing \textit{encoder}, model mampu menangkap pola visual dari gambar serta makna linguistik dari teks yang saling melengkapi dalam proses klasifikasi.

    \item Integrasi representasi visual dan tekstual melalui pendekatan multimodal terbukti memberikan peningkatan performa yang signifikan dibandingkan model unimodal. Hal ini ditunjukkan oleh peningkatan nilai \textit{F1-score} dari 91,5\% pada model unimodal gambar dan 92,0\% pada model unimodal teks menjadi 96,1\% pada model multimodal. Peningkatan ini menunjukkan bahwa kombinasi kedua modalitas mampu memberikan pemahaman konteks yang lebih baik, terutama pada data meme yang mengandung makna implisit. Dengan demikian, pendekatan multimodal lebih efektif dalam menangani permasalahan klasifikasi yang melibatkan hubungan kompleks antara gambar dan teks.

    \item Evaluasi kinerja model menggunakan \textit{confusion matrix} menunjukkan bahwa model multimodal memiliki performa yang baik dalam mendeteksi meme \textit{self-harm}, dengan \textit{precision} yang tinggi serta \textit{recall} yang lebih baik dibandingkan model unimodal. Hal ini menunjukkan bahwa model mampu mengurangi kesalahan klasifikasi, khususnya \textit{false negative} yang krusial dalam deteksi konten berbahaya. Selain itu, hasil eksperimen juga menunjukkan bahwa penggunaan arsitektur yang lebih kompleks, yaitu \textit{Transformer Encoder}, tidak memberikan peningkatan performa yang signifikan dibandingkan metode \textit{simple fusion}, sehingga pendekatan yang lebih sederhana tetap efektif dalam menangkap hubungan antar modalitas.

    \item Evaluasi tambahan menggunakan \textit{annotator} manusia menunjukkan bahwa meskipun model memiliki performa yang baik secara kuantitatif, masih terdapat tantangan dalam memahami data yang bersifat ambigu dan di luar distribusi data. Hasil analisis menunjukkan bahwa model cenderung mengalami \textit{false negative} pada meme yang mengandung makna implisit, seperti penggunaan permainan kata, sarkasme, dan humor gelap, serta \textit{false positive} pada konten yang menampilkan ekspresi emosional tanpa indikasi \textit{self-harm}. Hal ini mengindikasikan bahwa model masih memiliki keterbatasan dalam memahami hubungan multimodal secara kontekstual, terutama ketika terjadi konflik antara teks dan visual.
\end{enumerate}

Secara keseluruhan, penelitian ini menunjukkan bahwa pendekatan multimodal memberikan peningkatan performa yang signifikan dalam klasifikasi meme \textit{self-harm}, namun masih memerlukan pengembangan lebih lanjut dalam menangani kasus-kasus ambigu dan kompleks. Oleh karena itu, penambahan \textit{dataset} dengan variasi konteks yang lebih luas serta peningkatan kemampuan model dalam memahami makna implisit menjadi arah pengembangan yang penting untuk penelitian selanjutnya.


\section{Saran} \label{V.Saran}
Berdasarkan penelitian yang telah dilakukan, maka saran yang dapat diajukan adalah sebagai berikut:

\begin{enumerate}
    \item Penelitian selanjutnya disarankan untuk melibatkan ahli di bidang psikologi atau kesehatan mental dalam proses verifikasi hasil anotasi data. Meskipun pendekatan \textit{ensemble pseudo-labeling} yang digunakan dalam penelitian ini mampu menghasilkan anotasi dengan tingkat kesepakatan yang tinggi, validasi oleh ahli tetap diperlukan untuk memastikan akurasi dan sensitivitas dalam mengidentifikasi konten \textit{self-harm}, terutama pada kasus yang bersifat implisit atau ambigu. Hal ini menjadi penting mengingat hasil evaluasi menunjukkan bahwa interpretasi terhadap konten \textit{self-harm} tidak selalu bersifat objektif dan dapat dipengaruhi oleh konteks serta perspektif individu. Dengan demikian, keterlibatan ahli diharapkan dapat meningkatkan kualitas label \textit{dataset}, sehingga model yang dihasilkan menjadi lebih reliabel dan sesuai dengan konteks klinis maupun sosial yang sebenarnya.

    \item Penelitian selanjutnya juga disarankan untuk memperkaya \textit{dataset} dengan variasi data yang lebih beragam, khususnya yang mengandung unsur ambiguitas seperti \textit{wordplay}, sarkasme, humor gelap, serta konflik antara teks dan visual. Hal ini didasarkan pada temuan bahwa model masih mengalami kesulitan dalam mendeteksi konten \textit{self-harm} pada kasus-kasus tersebut, yang tercermin dari munculnya \textit{false negative} pada data yang bersifat implisit. Dengan memperluas cakupan \textit{dataset}, diharapkan model dapat belajar memahami pola-pola kompleks dan meningkatkan kemampuan generalisasi terhadap data di luar distribusi pelatihan (\textit{out-of-distribution}).

    \item Penelitian selanjutnya juga dapat mengembangkan sistem deteksi yang lebih aplikatif, seperti integrasi model ke dalam sistem monitoring media sosial secara \textit{real-time}. Pengembangan ini dapat dimanfaatkan sebagai langkah preventif dalam mengidentifikasi dan menangani konten \textit{self-harm} secara lebih dini. Namun, implementasi sistem tersebut perlu mempertimbangkan aspek etika, privasi, serta potensi dampak sosial, sehingga penggunaannya tetap bertanggung jawab dan tidak menimbulkan bias atau kesalahan interpretasi yang merugikan.
\end{enumerate}
